\chapter{热力学的基本概念与基本规律}

\begin{gather*}
    \begin{cases}
        \text{传统热力学：准静态过程} \\
        \text{非平衡态热力学（线性理论）：热力学流与热力学力之间有线性关系} \\
        \text{非平衡态热力学（非线性理论）}
    \end{cases}
\end{gather*}

热力学系统：大量微观粒子组成的宏观物体

\section{平衡态}

术语
\begin{itemize}
    \item 绝热壁、导热壁
    \item 刚性壁
    \item 热接触
    \item 孤立系
    \item 闭系、开系
\end{itemize}

\begin{gather*}
    \begin{cases}
        \text{平衡态：在没有外界影响下，物体的各部分性质长时间内不发生任何变化的状态}
        \\
        \text{稳恒态：存在外界影响下，物体的各部分性质长时间内不发生任何变化的状态}
    \end{cases}
\end{gather*}

\begin{itemize}
    \item 动态平衡
    \item 在一定条件下，初始不处于平衡态的系统，经过弛豫时间，必将趋于平衡态
    \begin{itemize}
        \item 孤立系
        \item 维持不变的外界条件
        \begin{itemize}
            \item 恒等温度：热库
            \item 恒定压强
            \item 开系：粒子库
        \end{itemize}
    \end{itemize}
\end{itemize}

平衡态的描写：宏观状态变量
\begin{itemize}
    \item 几何变量
    \item 力学变量
    \item 电磁变量
    \item 化学变量
    \item 热力学变量：温度
\end{itemize}
对于均匀系，分为广延量和强度量

\begin{itemize}
    \item 准静态过程：在过程中系统始终处于平衡态，外界变化的特征时间远大于弛豫时间
    \item 可逆过程：每一步都可以在相反方向进行而不引起外界其他变化的过程
\end{itemize}

\section{物态方程}

温度与独立状态变量$(x_1, x_2, \ldots, x_n)$之间的关系
\begin{gather*}
    f(T; x_1, x_2, \ldots, x_n) = 0
\end{gather*}

对于$p-V-T$系统
\begin{itemize}
    \item 膨胀系数
    \begin{gather*}
        \alpha = \frac{1}{V} \left( \frac{\partial V}{\partial T} \right)_p
    \end{gather*}
    \item 压强系数
    \begin{gather*}
        \beta = \frac{1}{p} \left( \frac{\partial p}{\partial T} \right)_V
    \end{gather*}
    \item 等温压缩系数
    \begin{gather*}
        \kappa_T = -\frac{1}{V} \left( \frac{\partial V}{\partial p} \right)_T
    \end{gather*}
\end{itemize}
满足
\begin{gather*}
    \alpha = \beta \kappa_T p
\end{gather*}

\section{热力学第一定律}

\begin{gather*}
    \d U = \d Q + \d W
\end{gather*}
$Q$为吸热，$W$为外界对系统做功

\subsection{温度}

\begin{itemize}
    \item 热平衡：热接触后状态不发生变化
    \item 热平衡定律$\implies$一切互为热平衡的物体的温度相等
\end{itemize}

\subsection{内能}

初末态不同绝热过程的功相等，因此定义态函数内能
\begin{gather*}
    \Delta U = W_{ad}
\end{gather*}

\subsection{功}

外界对系统做功
\begin{itemize}
    \item 准静态过程的功
    \begin{itemize}
        \item 流体体积变化
        \begin{gather*}
            \d W = -p \d V
        \end{gather*}
        \item 表面膜面积变化
        \begin{gather*}
            \d W = \sigma \d A
        \end{gather*}
        \item 细弹性丝长度变化
        \begin{gather*}
            \d W = F \d L
        \end{gather*}
        \item 电介质极化
        \begin{gather*}
            \d W = V \vec{\mathscr{E}} \cdot \d \vec{\mathscr{D}} = V \d \left(\frac{1}{2} \varepsilon_0 \mathscr{E}^2 \right) + V \vec{\mathscr{E}} \cdot \d \vec{\mathscr{P}}
        \end{gather*}
        第一项为真空电场能量变化，第二项为极化功
        \item 磁介质磁化
        \begin{gather*}
            \d W = V \vec{\mathscr{H}} \cdot \d \vec{\mathscr{B}} = V \d \left(\frac{1}{2\mu_0} \mathscr{B}^2 \right) + \mu_0 V \vec{\mathscr{H}} \cdot \d \vec{\mathscr{M}}
        \end{gather*}
    \end{itemize}
    \item 非准静态过程的功
    \begin{itemize}
        \item 等容过程
        \begin{gather*}
            W = 0
        \end{gather*}
        \item 等压过程
        \begin{gather*}
            W = -p_{ex} \Delta V
        \end{gather*}
        $p_{ex}$为外界压强
    \end{itemize}
\end{itemize}

\subsection{热容}

过程$y$中吸收的热量为$Q_y$，定义热容
\begin{gather*}
    C_y = \frac{\d Q_y}{\d T}
\end{gather*}

\section{热力学第二定律}

\begin{itemize}
    \item Clausius表述
    \begin{gather*}
        \text{不可能有一种循环过程，其唯一结果是将热量从低温物体传递到高温物体}
    \end{gather*}
    \item Kelvin表述
    \begin{gather*}
        \text{不可能有一种循环过程，其唯一结果是从单一热源吸收热量并将其全部转化为功}
    \end{gather*}
\end{itemize}

\begin{itemize}
    \item Carnot定理：工作于两个热库之间的热机，可逆热机效率$\eta$最大
    \item Clausius不等式：从$n$个热库吸收热量$Q_i$，温度$T_i$，对任意循环过程
    \begin{gather*}
        \sum_{i=1}^{n} \frac{Q_i}{T_i} \leq 0
    \end{gather*}
    连续化，对于任意循环过程，
    \begin{gather*}
        \int \frac{\d Q}{T} \leq 0
    \end{gather*}
    \item 熵：
    \begin{gather*}
        S_B - S_A = \int_A^B \frac{\d Q_{rev}}{T}
    \end{gather*}
    积分路径为任一可逆过程
    \begin{gather*}
        \d S \geq \frac{\d Q}{T}
    \end{gather*}
    $T$代表热源温度
    \item 熵增加原理：系统的熵在绝热过程中永不减少，在可逆绝热过程中不变，在不可逆绝热过程中增加
    \item 最大功原理：可逆过程输出的功最大
\end{itemize}